\scanpage{27}
\mspage{5}

$P$, et un diviseur \inserted{relatif positif} $\Theta$ sur $P$, [tel que $\deg\Theta^g/g!=1$] et que la polarisation $\delta$ soit égale à celle déduite de la polarisation de $P$ par le diviseur $\Theta$. Etant donné \deleted{\unsure{Et}} $P'$ et $\Theta'$ ayant les \unsure{$=^{s}$|mêmes} vertus que $P$ et $\Theta$, il existe un unique $A$-isom de $P$ sur $P'$, transformant $\Theta$ en $\Theta'$ [en d'autres termes \unsure{toute|seule} \deleted{\unsure{autre}} section de $A$ qui invarie $\Theta$ est la section nulle].

b) Si on a un torseur $P$ sous $A$, muni d'un diviseur relatif positif $\Theta$ \deleted{de degré de polarisation 1,}\inserted{tel que $\deg\Theta^g/g!=1$}, alors $\Theta$ est ample rel $/S$ et définit \deleted{\unsure{donc}} une polarisation de $P/S$ donc \deleted{degré} de $A/S$. Cette dernière est de degré de polarisation 1.

c) Sous les conditions de b), considérons l'hom de composition $A\times P\longrightarrow P$ ; et soit $\Theta'$ l'image inverse de $\Theta$ dans $A\times P$. Alors $\Theta$ est un diviseur positif relatif sur $A\times P\inserted{/P}$, et définit donc un \deleted{\unsure{f}} faisceau inversible sur $A\times P$, donc un hom $P\mathbin{\unsure{-|\to}}\underline{\mathrm{Pic}}_{A/S}$. Ce dernier est un isom de $P$ sur $\underline{\mathrm{Pic}}^\delta_{A/S}$, où $\delta$ est la polarisation de $A/S$ définie par \deleted{$\unsure{\delta}$} $\Theta$ sur $P$.

d) Supposons que $A=\underline{\mathrm{Alb}}^0_{X/S}$, où $X/S$ est une courbe propre et lisse relative de genre $g$. \deleted{Alors} Considérons la polarisation canonique $\delta$ de $A$ \deleted{\unsure{de $B$, d'où un iso}}\inserted{\unsure{définie|défini} par le diviseur diagonal} sur $X\times_S X$. Cette polarisation est de degré 1, et $P^\delta$ est canoniquement isomorphe à \inserted{comme $A$-torseur} $\underline{\mathrm{Pic}}^{g-1}_{X/S}=\underline{\mathrm{Alb}}^{g-1}_{X/S}$, est isomorphe\inserted{ à} \unsure{sous|sur}
% slanted

\margin{question de signe!}
% archive sheet number: 26
